% Frozen isolated research proof, v2.  This file is not a live-book input.
% It has not been submitted to production Rethlas verification.

We first isolate the only ordinal fact that is needed.  Every ordinal below
\(\omega _1\) is countable, and the supremum of a countable family of
countable ordinals is again countable.  Hence, if
\(\beta_n<\omega _1\) for all \(n\), then
\begin{equation}
 \label{eq:p072-countable-supremum}
                  \sup_{n\geq1}\beta_n<\omega _1.
\end{equation}

At the limit ordinal \(\omega _1\), clause~(c) of
\cite[Lemma~8.1.4]{CLD25} is exactly
\eqref{eq:p072-omega-one-local-union}.  This is a local, or
sheaf-theoretic, union rather than the assertion that a section on an
arbitrary open set belongs globally to one fixed earlier stage.  In stalk
notation it says
\begin{equation}
 \label{eq:p072-stalkwise-union}
  \bigl(\operatorname{Psh}_{X,\omega _1}\bigr)_x
  =
  \bigcup_{\alpha<\omega _1}
  \bigl(\operatorname{Psh}_{X,\alpha}\bigr)_x
  \qquad(x\in X).
\end{equation}

Fix \(x\in U\).  The topology underlying an affine \(G\)-tropical space is
Hausdorff and locally compact by \cite[\S3.1.1]{CLD25}.  We may therefore
choose a compact set \(K\) such that
\begin{equation}
 \label{eq:p072-compact-neighbourhood}
                    x\in K^\circ\subseteq K\subseteq U.
\end{equation}
Fix \(n\).  Since \(\omega _1\) is a limit ordinal and
\(u_n\in\operatorname{Psh}_{X,\omega _1}(U)\), the definition of a limit
stage says that every \(y\in K\) has an open neighbourhood \(W_{n,y}\)
and an ordinal \(\beta_{n,y}<\omega _1\) for which
\begin{equation}
 \label{eq:p072-local-rank}
     u_n|_{W_{n,y}}\in
     \operatorname{Psh}_{X,\beta_{n,y}}(W_{n,y}).
\end{equation}
After replacing \(W_{n,y}\) by \(W_{n,y}\cap U\), all these neighbourhoods
may be taken inside \(U\).

Compactness of \(K\) gives points \(y_{n,1},\ldots,y_{n,r_n}\) such that
the corresponding \(W_{n,y_{n,j}}\) cover \(K\).  Put
\begin{equation}
 \label{eq:p072-row-rank}
       \beta_n=\max_{1\leq j\leq r_n}\beta_{n,y_{n,j}},
       \qquad
       W_n=\bigcup_{j=1}^{r_n}W_{n,y_{n,j}}.
\end{equation}
The filtration is increasing.  Thus every restriction in
\eqref{eq:p072-local-rank} belongs to
\(\operatorname{Psh}_{X,\beta_n}\).  These restrictions agree on overlaps,
because they are restrictions of the one function \(u_n\).  Since
\(\operatorname{Psh}_{X,\beta_n}\) is a sheaf, they glue to give
\begin{equation}
 \label{eq:p072-row-common-core}
              u_n|_{W_n}\in
              \operatorname{Psh}_{X,\beta_n}(W_n).
\end{equation}
Notice that \(K\subseteq W_n\) for every \(n\).  Therefore, on the one fixed
open neighbourhood
\begin{equation}
 \label{eq:p072-fixed-open-core}
                              V=K^\circ
\end{equation}
of \(x\), all the rows in \eqref{eq:p072-row-common-core} are defined.

Set
\begin{equation}
 \label{eq:p072-uniform-rank}
                         \alpha=\sup_{n\geq1}\beta_n.
\end{equation}
Equation \eqref{eq:p072-countable-supremum} gives
\(\alpha+1<\omega _1\).  Monotonicity of the filtration and restriction of
\eqref{eq:p072-row-common-core} yield
\begin{equation}
 \label{eq:p072-all-rows-one-stage}
       u_n|_V\in\operatorname{Psh}_{X,\alpha}(V)\qquad(n\geq1).
\end{equation}
The sequence in \eqref{eq:p072-all-rows-one-stage} decreases pointwise to
\(u|_V\).  The definition of the successor stage now gives
\begin{equation}
 \label{eq:p072-limit-next-stage}
              u|_V\in\operatorname{Psh}_{X,\alpha+1}(V).
\end{equation}
Since \(\alpha+1<\omega _1\), this proves
\eqref{eq:p072-common-countable-rank}.  As \(x\) was arbitrary, it also
gives an open cover of \(U\) on which \(u\) belongs to
\(\operatorname{Psh}_{X,\omega _1}\).  The latter is a sheaf, so
\(u\in\operatorname{Psh}_{X,\omega _1}(U)\).  Thus
\(\operatorname{Psh}_{X,\omega _1}\) is stable under decreasing sequential
limits.

The sheaf \(\operatorname{Psh}_{X,\omega _1}\) contains the smooth psh
generators.  We have just proved that it is stable under the operation in
the definition of \(\operatorname{Psh}_X\).  The defining minimality of
\(\operatorname{Psh}_X\) therefore gives
\begin{equation}
 \label{eq:p072-minimality-first-inclusion}
                 \operatorname{Psh}_X
                 \subseteq\operatorname{Psh}_{X,\omega _1}.
\end{equation}
Conversely, transfinite induction in
\cite[Lemma~8.1.4]{CLD25} gives
\(\operatorname{Psh}_{X,\gamma}\subseteq\operatorname{Psh}_X\) at every
ordinal \(\gamma\).  Taking \(\gamma=\omega _1\) and combining with
\eqref{eq:p072-minimality-first-inclusion} proves
\eqref{eq:p072-omega-one-psh}.

For the continuous hierarchy, CLD prove
\begin{equation}
 \label{eq:p072-cld-continuous-identification}
 \operatorname{CPsh}_{X,\gamma}
 =C_X^0\cap\operatorname{Psh}_{X,\gamma}
 \quad\text{for every ordinal }\gamma,
 \qquad
 \operatorname{CPsh}_X=C_X^0\cap\operatorname{Psh}_X
\end{equation}
\cite[Proposition~8.1.11]{CLD25}.  Taking \(\gamma=\omega _1\) in
\eqref{eq:p072-cld-continuous-identification} and using
\eqref{eq:p072-omega-one-psh} gives
\[
 \operatorname{CPsh}_{X,\omega _1}
 =C_X^0\cap\operatorname{Psh}_{X,\omega _1}
 =C_X^0\cap\operatorname{Psh}_X
 =\operatorname{CPsh}_X,
\]
which is \eqref{eq:p072-omega-one-cpsh}.

The argument proves the universal first-uncountable-stage bound only.  It
does not give a fixed countable ordinal, and it does not show first-stage
collapse: the countable ranks \(\beta_n\) in
\eqref{eq:p072-row-rank} may be unbounded in \(\omega _1\) as the input and
the point vary.
